% Bài 19 — SGK Toán 9 KNTT Tập 2, trang 10–17.
% Nguồn direct image theo issue #3:
% https://cdn3.olm.vn/upload/thuvienso/4491669493-page-10-1771844538088.png
% https://cdn3.olm.vn/upload/thuvienso/4491669493-page-11-1771844538428.png
% https://cdn3.olm.vn/upload/thuvienso/4491669493-page-12-1771844538768.png
% https://cdn3.olm.vn/upload/thuvienso/4491669493-page-13-1771844539117.png
% https://cdn3.olm.vn/upload/thuvienso/4491669493-page-14-1771844539459.png
% https://cdn3.olm.vn/upload/thuvienso/4491669493-page-15-1771844539798.png
% https://cdn3.olm.vn/upload/thuvienso/4491669493-page-16-1771844540117.png
% https://cdn3.olm.vn/upload/thuvienso/4491669493-page-17-1771844540426.png
%
% SBT Toán 9 KNTT Tập 2 — Bài 19, bài tập 6.9–6.16.
% Authority: Project Source "SBT TOÁN 9 TẬP 2.pdf", PDF pages 10–11.
% SHA-256: 5ff01fd213d1adc75f9c54b4408d6a9642a4d8038e685040efef3311a196196c
% Không lấy phần LỜI GIẢI – HƯỚNG DẪN – ĐÁP SỐ.

\section{Phương trình bậc hai một ẩn}

\begin{lesson}
    \subsection{Kiến thức trọng tâm}

    Trên một mảnh đất hình chữ nhật có kích thước $28$ m $\times$ $16$ m, người ta dự định làm một bể bơi có đường đi xung quanh (H.6.9). Hỏi bề rộng của đường đi là bao nhiêu để diện tích của bể bơi là $288\,\text{m}^2$?

    % TODO(HINH-NGUON): H.6.9 là ảnh bể bơi; bổ sung asset/crop đúng từ SGK trang 10, không redraw xấp xỉ bằng TikZ.

    \subsubsection{Định nghĩa phương trình bậc hai một ẩn}

    \textbf{Nhận biết phương trình bậc hai một ẩn}

    Xét bài toán trong tình huống mở đầu.

    \textbf{HĐ1.} Gọi $x$ (m) là bề rộng của mặt đường ($0<x<8$). Tính chiều dài và chiều rộng của bể bơi theo $x$.

    \answerline
    \answerline

    \textbf{HĐ2.} Dựa vào kết quả HĐ1, tính diện tích của bể bơi theo $x$.

    \answerline
    \answerline

    \textbf{HĐ3.} Sử dụng giả thiết và kết quả HĐ2, hãy viết phương trình để tìm $x$.

    \answerline
    \answerline

    \begin{keyconcept}
        Phương trình bậc hai một ẩn (nói gọn là phương trình bậc hai) là phương trình có dạng
        \[
            ax^2+bx+c=0,
        \]
        trong đó $x$ là ẩn; $a$, $b$, $c$ là những số cho trước gọi là hệ số và $a\ne0$.
    \end{keyconcept}

    \textbf{Luyện tập 1.} Trong các phương trình sau, những phương trình nào là phương trình bậc hai ẩn $x$? Chỉ rõ các hệ số $a$, $b$, $c$ của mỗi phương trình bậc hai đó.
    \begin{enumerate}[label=\alph*)]
        \item $x^2+5=0$;
        \item $2x^2+7x=0$;
        \item $\dfrac{x^2-2x+5}{x}=0$;
        \item $0{,}5x^2=0$.
    \end{enumerate}

    \answerline
    \answerline

    \textbf{Tranh luận.} Phương trình (ẩn $x$) $mx^2+2x+1=0$ ($m$ là một số cho trước) là một phương trình bậc hai với $a=m$, $b=2$, $c=1$.

    Ý kiến của em thế nào?

    \answerline
    \answerline

    \subsubsection{Cách giải phương trình bậc hai một ẩn có dạng đặc biệt}

    \textbf{Cách giải phương trình bậc hai một ẩn dạng khuyết}

    Giải một phương trình bậc hai là tìm tất cả các nghiệm của nó. Dưới đây, thông qua một số ví dụ đơn giản, ta trình bày cách giải một số phương trình bậc hai dạng $ax^2+bx+c=0$ ($a\ne0$), mà khuyết số hạng bậc nhất (tức là $b=0$) hoặc khuyết số hạng tự do (tức là $c=0$), bằng phương pháp đặt nhân tử chung đưa về dạng tích hoặc dùng hằng đẳng thức để đưa vế trái về một bình phương.

    \begin{keyconcept}
        \begin{itemize}
            \item Nếu $A\cdot B=0$ thì $A=0$ hoặc $B=0$.
            \item Nếu $A^2=B$ ($B\ge0$) thì $A=\sqrt{B}$ hoặc $A=-\sqrt{B}$.
        \end{itemize}
    \end{keyconcept}

    \textbf{Luyện tập 2.} Giải các phương trình sau:
    \begin{enumerate}[label=\alph*)]
        \item $2x^2+6x=0$;
        \item $5x^2+11x=0$.
    \end{enumerate}

    \answerline
    \answerline

    \textbf{Luyện tập 3.} Giải các phương trình sau:
    \begin{enumerate}[label=\alph*)]
        \item $x^2-25=0$;
        \item $(x+3)^2=5$.
    \end{enumerate}

    \answerline
    \answerline

    \textbf{Chú ý.} Để giải phương trình bậc hai dạng $x^2+bx=c$, ta có thể cộng thêm vào hai vế của phương trình với cùng một số thích hợp để vế trái có thể biến đổi thành một bình phương. Từ đó có thể giải phương trình đã cho.

    \textbf{Luyện tập 4.} Cho phương trình $x^2+6x=1$. Hãy cộng vào cả hai vế của phương trình với cùng một số thích hợp để được một phương trình mà vế trái có thể biến đổi thành một bình phương. Từ đó, hãy giải phương trình đã cho.

    \answerline
    \answerline
    \answerline

    \subsubsection{Công thức nghiệm của phương trình bậc hai}

    \textbf{Xây dựng công thức nghiệm của phương trình bậc hai}

    \textbf{HĐ4.} Thực hiện lần lượt các bước sau để giải phương trình:
    \[
        2x^2-8x+3=0.
    \]
    \begin{enumerate}[label=\alph*)]
        \item Chuyển hạng tử tự do sang vế phải.
        \item Chia cả hai vế của phương trình cho hệ số của $x^2$.
        \item Thêm vào hai vế của phương trình nhận được ở câu b với cùng một số để vế trái có thể biến đổi thành một bình phương. Từ đó tìm nghiệm $x$.
    \end{enumerate}

    \answerline
    \answerline
    \answerline

    \textbf{Cách giải phương trình bậc hai}

    Tương tự HĐ4, để giải phương trình bậc hai $ax^2+bx+c=0$ ($a\ne0$) trong trường hợp tổng quát, ta làm như sau:
    \begin{itemize}
        \item Chuyển hạng tử tự do $c$ sang vế phải: $ax^2+bx=-c$.
        \item Chia cả hai vế của phương trình cho hệ số $a$ của $x^2$: $x^2+\dfrac{b}{a}x=-\dfrac{c}{a}$.
        \item Cộng vào hai vế của phương trình nhận được với $\dfrac{b^2}{4a^2}$ để vế trái có thể biến đổi thành bình phương của một biểu thức:
        \[
            x^2+\dfrac{b}{a}x+\dfrac{b^2}{4a^2}=-\dfrac{c}{a}+\dfrac{b^2}{4a^2},
        \]
        hay
        \[
            \left(x+\dfrac{b}{2a}\right)^2=\dfrac{b^2-4ac}{4a^2}.
        \]
    \end{itemize}

    Kí hiệu $\Delta=b^2-4ac$ và gọi là biệt thức của phương trình ($\Delta$ đọc là ``đenta''). Khi đó, ta có thể viết lại phương trình cuối dưới dạng
    \[
        \left(x+\dfrac{b}{2a}\right)^2=\dfrac{\Delta}{4a^2}.
    \]

    Từ đây, ta có kết quả sau:

    \begin{keyconcept}
        Xét phương trình bậc hai một ẩn $ax^2+bx+c=0$ ($a\ne0$).

        Tính biệt thức $\Delta=b^2-4ac$.
        \begin{itemize}
            \item Nếu $\Delta>0$ thì phương trình có hai nghiệm phân biệt:
            \[
                x_1=\dfrac{-b+\sqrt{\Delta}}{2a};\qquad
                x_2=\dfrac{-b-\sqrt{\Delta}}{2a}.
            \]
            \item Nếu $\Delta=0$ thì phương trình có nghiệm kép
            \[
                x_1=x_2=-\dfrac{b}{2a}.
            \]
            \item Nếu $\Delta<0$ thì phương trình vô nghiệm.
        \end{itemize}
    \end{keyconcept}

    \textbf{Luyện tập 5.} Áp dụng công thức nghiệm, giải các phương trình sau:
    \begin{enumerate}[label=\alph*)]
        \item $2x^2-5x+1=0$;
        \item $x^2+8x+16=0$;
        \item $x^2-x+1=0$.
    \end{enumerate}

    \answerline
    \answerline
    \answerline

    \textbf{Thử thách nhỏ.} Có thể nói gì về nghiệm của phương trình bậc hai $ax^2+bx+c=0$ nếu $a$ và $c$ trái dấu?

    Em hãy trả lời câu hỏi của anh Pi.

    \answerline
    \answerline

    \textbf{Chú ý.} Xét phương trình bậc hai $ax^2+bx+c=0$ ($a\ne0$), với $b=2b'$ và $\Delta'=b'^2-ac$.
    \begin{itemize}
        \item Nếu $\Delta'>0$ thì phương trình có hai nghiệm phân biệt:
        \[
            x_1=\dfrac{-b'+\sqrt{\Delta'}}{a};\qquad
            x_2=\dfrac{-b'-\sqrt{\Delta'}}{a}.
        \]
        \item Nếu $\Delta'=0$ thì phương trình có nghiệm kép
        \[
            x_1=x_2=-\dfrac{b'}{a}.
        \]
        \item Nếu $\Delta'<0$ thì phương trình vô nghiệm.
    \end{itemize}

    Các công thức ở trên gọi là công thức nghiệm thu gọn.

    \textbf{Luyện tập 6.} Xác định $a$, $b'$, $c$ rồi dùng công thức nghiệm thu gọn giải các phương trình sau:
    \begin{enumerate}[label=\alph*)]
        \item $3x^2+8x-3=0$;
        \item $x^2+6\sqrt{2}x+2=0$.
    \end{enumerate}

    \answerline
    \answerline

    \textbf{Vận dụng.} Giải bài toán trong tình huống mở đầu.

    \answerline
    \answerline
    \answerline

    \subsubsection{Tìm nghiệm của phương trình bậc hai bằng máy tính cầm tay}

    Sử dụng máy tính cầm tay, ta có thể dễ dàng tìm nghiệm của các phương trình bậc hai một ẩn.

    \textbf{Luyện tập 7.} Sử dụng máy tính cầm tay, tìm nghiệm của các phương trình sau:
    \begin{enumerate}[label=\alph*)]
        \item $5x^2+2\sqrt{10}x+2=0$;
        \item $3x^2-5x+7=0$;
        \item $4x^2-11x+1=0$.
    \end{enumerate}

    \answerline
    \answerline
\end{lesson}

\begin{exercises}
    \subsection{Bài tập}

    \begin{questions}[resume=SGK]
        \item Đưa các phương trình sau về dạng $ax^2+bx+c=0$ và xác định các hệ số $a$, $b$, $c$ của phương trình đó.
        \begin{questions}
            \item $3x^2+2x-1=x^2-x$;
            \item $(2x+1)^2=x^2+1$.
        \end{questions}
        \answerline
        \answerline

        \item Giải các phương trình sau:
        \begin{questions}
            \item $2x^2+\dfrac{1}{3}x=0$;
            \item $(3x+2)^2=5$.
        \end{questions}
        \answerline
        \answerline

        \item Không cần giải phương trình, hãy xác định các hệ số $a$, $b$, $c$, tính biệt thức $\Delta$ và xác định số nghiệm của mỗi phương trình sau:
        \begin{questions}
            \item $11x^2+13x-1=0$;
            \item $9x^2+42x+49=0$;
            \item $x^2-2x+3=0$.
        \end{questions}
        \answerline
        \answerline

        \item Dùng công thức nghiệm của phương trình bậc hai, giải các phương trình sau:
        \begin{questions}
            \item $x^2-2\sqrt{5}x+2=0$;
            \item $4x^2+28x+49=0$;
            \item $3x^2-3\sqrt{2}x+1=0$.
        \end{questions}
        \answerline
        \answerline

        \item Sử dụng máy tính cầm tay, tìm nghiệm của các phương trình sau:
        \begin{questions}
            \item $0{,}1x^2+2{,}5x-0{,}2=0$;
            \item $0{,}01x^2-0{,}05x+0{,}0625=0$;
            \item $1{,}2x^2+0{,}75x+2{,}5=0$.
        \end{questions}
        \answerline
        \answerline

        \item Độ cao $h$ (mét) so với mặt đất của một vật được phóng thẳng đứng lên trên từ mặt đất với vận tốc ban đầu $v_0=19{,}6$ m/s cho bởi công thức $h=19{,}6t-4{,}9t^2$, ở đó $t$ là thời gian kể từ khi phóng (giây) (theo \textit{Vật lí đại cương}, NXB Giáo dục Việt Nam, 2016). Hỏi sau bao lâu kể từ khi phóng, vật sẽ rơi trở lại mặt đất?

        \answerline
        \answerline

        \item Kích thước màn hình ti vi hình chữ nhật được xác định bằng độ dài đường chéo. Ti vi truyền thống có định dạng $4:3$, nghĩa là tỉ lệ giữa chiều dài và chiều rộng của màn hình là $4:3$. Hỏi diện tích của màn hình ti vi truyền thống $37$ in là bao nhiêu? Diện tích của màn hình ti vi LCD $37$ in có định dạng $16:9$ là bao nhiêu? Màn hình ti vi nào có diện tích lớn hơn? Ở đây, các diện tích của màn hình được tính bằng inch vuông.

        % TODO(HINH-NGUON): bổ sung đúng hai hình ti vi từ SGK trang 17 nếu cần asset; không redraw ảnh minh họa bằng TikZ.

        \answerline
        \answerline
        \answerline

        \item Một mảnh vườn hình chữ nhật có chiều rộng ngắn hơn chiều dài $6$ m và có diện tích là $140\,\text{m}^2$. Tính các kích thước của mảnh vườn đó.

        \answerline
        \answerline
    \end{questions}
\end{exercises}

\begin{practice}
    \subsection{Luyện tập}

    \begin{questions}[resume=SBT]
        \item Giải các phương trình sau bằng cách đưa về dạng tích:
        \begin{questions}
            \item $x^2+5x=0$;
            \item $x^2-16=0$;
            \item $x^2-10x+25=0$;
            \item $x^2+8x+12=0$.
        \end{questions}
        \answerline
        \answerline

        \item Giải các phương trình sau:
        \begin{questions}
            \item $(2x+1)^2=3$;
            \item $(2-3x)^2=5$.
        \end{questions}
        \answerline
        \answerline

        \item Sử dụng công thức nghiệm hoặc công thức nghiệm thu gọn, giải các phương trình bậc hai sau:
        \begin{questions}
            \item $x^2+2x-5=0$;
            \item $4x^2-4\sqrt{3}x+3=0$;
            \item $x^2-6\sqrt{5}x+7=0$.
        \end{questions}
        \answerline
        \answerline

        \item Sử dụng máy tính cầm tay, tìm nghiệm của các phương trình sau:
        \begin{questions}
            \item $2x^2+\sqrt{11}x-1=0$;
            \item $\dfrac{1}{2}x^2+\dfrac{5}{3}x+\dfrac{50}{9}=0$;
            \item $\sqrt{2}x^2-(1+\sqrt{5})x+11=0$.
        \end{questions}
        \answerline
        \answerline

        \item Tùy theo các giá trị của $m$, hãy giải phương trình ẩn $x$ sau: $(2x-1)^2=m$.

        \answerline
        \answerline

        \item Cho phương trình (ẩn $x$): $x^2+4(m+1)x+4m^2-3=0$.
        \begin{questions}
            \item Tính biệt thức $\Delta'$.
            \item Tìm điều kiện của $m$ để phương trình:
            \begin{itemize}
                \item Có hai nghiệm phân biệt;
                \item Có nghiệm kép;
                \item Vô nghiệm.
            \end{itemize}
        \end{questions}
        \answerline
        \answerline
        \answerline

        \item Quỹ đạo chuyển động của một quả bóng được cho bởi công thức $y=1{,}5+x-0{,}098x^2$, trong đó $y$ (mét) là độ cao của quả bóng so với mặt đất và $x$ (mét) là khoảng cách theo phương ngang từ vị trí của quả bóng đến vị trí ném (xem hình bên). Tính khoảng cách theo phương ngang từ vị trí ném bóng đến vị trí quả bóng chạm đất.

        \begin{center}
        \begin{tikzpicture}[x=0.42cm,y=0.72cm]
            \draw[->] (0,0) -- (12.2,0) node[right] {$x$};
            \draw[->] (0,0) -- (0,5.0) node[above] {$y$};
            \draw[domain=0:11.55,samples=80,smooth,variable=\x]
                plot ({\x},{1.5+\x-0.098*\x*\x});
            \node[below left] at (0,0) {$O$};
            \node[left] at (0,1.5) {$1{,}5$};
            \node[above right] at (3.3,3.85) {$y=1{,}5+x-0{,}098x^2$};
        \end{tikzpicture}
        \end{center}

        \answerline
        \answerline

        \item Công ty sản xuất ván gỗ cần ước tính chiều dài tấm ván (tính bằng feet) có thể tạo ra được từ một khúc gỗ. Một trong những công thức được sử dụng phổ biến để ước tính chiều dài tấm ván là công thức Doyle:
        \[
            B=\dfrac{L}{16}(D^2-8D+16),
        \]
        trong đó $B$ là chiều dài tấm ván (feet), $D$ là đường kính (inch) và $L$ là chiều dài của khúc gỗ (feet).
        \begin{questions}
            \item Viết lại công thức Doyle cho các khúc gỗ dài $16$ feet.
            \item Tìm các nghiệm của phương trình bậc hai ẩn $D$ sau: $D^2-8D+16=0$.
        \end{questions}
        \answerline
        \answerline
    \end{questions}
\end{practice}
